\chapter{Conclusioni e Sviluppi Futuri}
Questo lavoro di tesi si è posto l'obiettivo di analizzare l'affidabilità dei Large Language Models nel dominio della Vulnerability Detection, approfondendo non solo la loro suscettibilità agli attacchi testuali ma mappando il processo di tale vulnerabilità attraverso la metodologia della Mechanistic Intepretability.

L'analisi condotta sulla baseline ha evidenziato come le architetture odierne, seppur addestrate specificamente per la programmazione, come le varianti coder, non possiedono ancora le capacità necessarie per operare in contesti di sicurezza senza la supervisione umana. Questa fragilità è stata confermata dai risultati degli attacchi Black-Box: modelli adottati in ambienti lavorativi come quelli della famiglia Qwen presentano un alto tasso di suscettibilità alle perturbazioni introdotte, superiore al 60\%. Gli altri modelli in particolare Llama-3.1-8B-Instruct hanno dimostrato di essere stati configurati in modo più resiliente rispetto a possibili attacchi.
Le metodologie di attacco presentate nello studio hanno portato alla luce diverse mancanze dei modelli, possibilmente fatali, registrando un tasso di successo molto alto. E'il caso della meccanica di Prompt Injection che registra un ASR del 100\% per il modello codellama.
L'inserimento di perturbazioni testuali, anche sintatticamente slegate dalla logica del codice sorgente analizzato, è sufficiente a sovrascrivere le direttive di sistema e a forzare una classificazione errata.

Il contributo principe di questa ricerca risiede tuttavia nell'esplorazione vettoriale del Residual Stream. Tramite il processo di estrazione implementato, è stato possibile superare la pura analisi comportamentale di input-output e visualizzare la topologia delle decisioni.
L'applicazione di metriche geometriche, quali la Cosine Similarity e la Distanza Euclidea, ha permesso la corretta interpretazione degli attacchi testuali nello spazio latente:
\begin{itemize}[noitemsep]
    \item Contrariamente a quanto si potesse ipotizzare, gli attacchi testuali non spingono il vettore della rappresentazione lungo la traiettoria della rappresentazione dei concetti di sicurezza e vulnerabilità della baseline. Le perturbazioni introdotte agiscono effettuando uno shift della traiettoria, deviando lo stato latente in dimensioni ortogonali. Il modello viene ingannato non perché riconosce il codice come sicuro, ma perché gli attacchi Black-Box distruggono la visione del task originario costringendo la rete ad operare in una zona grigia.
    \item Al contrario, l'attacco White-Box di Activation Steering si è dimostrato fedele alle ipotesi iniziali: nei layer di massima concentrazione computazionale lo spostamento avviene sull'asse del vettore di steering per poi allontanarsi solo nei livelli successivi del Transformer.
\end{itemize}
L'analisi congiunta tramti gli attacchi di Activation Steering, Logit Bias e strumenti di Logit Lens ha permesso di delineare le resistenze interne dei modelli stessi, evidenziando una tendenza al controllo dello spazio latente al costo di una debole resistenza nei layer finali.
Ciò evidenzia non solo la naturale evoluzione della computazione del modello, ma conferma le ipotesi di ripartizione dei compiti tra i layer del modello lasciando gli ultimi all'elaborazione sintattica della risposta mentre la vera decodifica logica della vulnerabilità avviene molto prima.

L'adozione di un modello con capacità di ragionamento differito come DeepSeek-R1 assieme alla modifica del prompt ha portato alla luce le vere potenzialità della Chain-of-Thought: il \textit{Thinking} del modello, sia indotto che nativo, ha portato ad una resistenza maggiore agli attacchi.

\section{Sviluppi Futuri}
Lo studio si propone di essere un punto di partenza per studi successivi: sebbene in letteratura si sostenga che il numero di parametri non influenzi la corretta identificazione di perturbazioni, risulta metodologicamente corretto avviare uno studio di ablazione adottando modelli con un numero di parametri più elevato rispetto a quelli presentati, oltre ad estendere la revisione a codici più lunghi per poter analizzare come l'attacco si sviluppa.
Inoltre, l'utilizzo di soli tre attacchi Black-Box non permette la completa visione dei parallelismi delle perturbazioni introdotte. Ampliare lo studio della cosine similarity a più attacchi di natura Black-Box aiuterebbe a mappare lo spettro su cui si collocano, permettendo in futuro di sviluppare tecniche di difesa o collaudo nella sicurezza dei modelli stessi.